\documentclass[10pt,twocolumn]{article}

\usepackage[a4paper,margin=1.65cm,columnsep=0.65cm]{geometry}
\usepackage[T1]{fontenc}
\usepackage[utf8]{inputenc}
\usepackage{lmodern}
\usepackage{microtype}
\usepackage{amsmath,amssymb,mathtools,bm}
\usepackage{booktabs,array}
\usepackage[dvipsnames]{xcolor}
\usepackage{enumitem}
\usepackage{listings}
\usepackage[hidelinks]{hyperref}
\usepackage[nameinlink,noabbrev]{cleveref}

\definecolor{codeblue}{RGB}{25,70,120}
\definecolor{codegray}{RGB}{245,246,248}
\definecolor{rulegray}{RGB}{205,210,216}
\hypersetup{
  colorlinks=true,
  linkcolor=codeblue,
  urlcolor=codeblue,
  citecolor=codeblue
}
\lstset{
  basicstyle=\ttfamily\footnotesize,
  backgroundcolor=\color{codegray},
  frame=single,
  rulecolor=\color{rulegray},
  columns=fullflexible,
  breaklines=true,
  keepspaces=true,
  showstringspaces=false,
  xleftmargin=2pt,
  xrightmargin=2pt
}
\setlist{nosep,leftmargin=*}
\setlength{\parindent}{1em}
\setlength{\parskip}{0pt}
\setlength{\abovedisplayskip}{5pt plus 2pt minus 2pt}
\setlength{\belowdisplayskip}{5pt plus 2pt minus 2pt}
\setlength{\textfloatsep}{8pt plus 2pt minus 2pt}
\setlength{\floatsep}{7pt plus 2pt minus 2pt}
\setlength{\intextsep}{7pt plus 2pt minus 2pt}
\emergencystretch=1.5em
\raggedbottom

\newcommand{\ii}{\mathrm{i}}
\newcommand{\e}{\mathrm{e}}
\newcommand{\dd}{\mathrm{d}}
\newcommand{\vloc}{V_{\mathrm{loc}}}
\newcommand{\vnl}{\widehat V_{\mathrm{nl}}}
\newcommand{\zion}{Z_{\mathrm{ion}}}
\newcommand{\gvec}{\mathbf G}
\newcommand{\rvec}{\mathbf r}
\newcommand{\Rvec}{\mathbf R}
\newcommand{\kvec}{\mathbf k}
\newcommand{\code}[1]{\texttt{#1}}
\DeclareMathOperator{\erf}{erf}
\DeclareMathOperator{\erfc}{erfc}

\title{\textbf{Pseudopotentials from First Principles to PyPWDFT}\\
\large A beginner's mathematical and code-level guide to the GTH implementation}
\author{Repository-specific technical note}
\date{\today}

\begin{document}
\maketitle

\begin{abstract}
This note explains why pseudopotentials are used, what information they
preserve, and how the separable Goedecker--Teter--Hutter (GTH)
pseudopotential in PyPWDFT is evaluated.  It begins with the all-electron
Kohn--Sham problem and develops the frozen-core, norm-conserving, local, and
non-local ideas before following the implementation from a parameter file to
the Hamiltonian and total energy.  The conventions and array shapes are those
of the current repository, so the final sections can be read beside
\path{pypwdft/gth.py}, \path{pypwdft/pypwdft.py}, and
\path{pypwdft/psystem.py}.
\end{abstract}

\section{The one-sentence idea}

A pseudopotential replaces the nuclear Coulomb field \emph{and} the chemically
inactive core electrons by a smoother effective ion that is designed to make
the remaining valence electrons behave like the all-electron valence
electrons outside the core region.

Three qualifications matter:
\begin{enumerate}
  \item It is not merely a softened Coulomb singularity.  It must reproduce
  angular-momentum-dependent valence scattering.
  \item It does not approximate every all-electron observable near a nucleus.
  Its purpose is accurate valence energetics and wavefunctions in the region
  relevant to bonding.
  \item It is part of the physical model.  A parameter set is tied to a
  valence partition and normally to an exchange--correlation (XC)
  approximation.
\end{enumerate}

\section{Where the difficulty comes from}

\subsection{The all-electron Kohn--Sham equation}

In atomic units (\(\hbar=m_e=e=4\pi\epsilon_0=1\)), a spin-unpolarized
Kohn--Sham orbital obeys
\begin{align}
\widehat H_{\mathrm{KS}}[\rho]
={}&-\frac{1}{2}\nabla^2+V_{\mathrm{nuc}}(\rvec)\nonumber\\
&+V_{\mathrm H}[\rho](\rvec)
+V_{\mathrm{xc}}[\rho](\rvec),\nonumber\\
\widehat H_{\mathrm{KS}}[\rho]\psi_n(\rvec)
={}&\epsilon_n\psi_n(\rvec).
\label{eq:ks-ae}
\end{align}
where
\begin{equation}
V_{\mathrm{nuc}}(\rvec)
=-\sum_a\frac{Z_a}{|\rvec-\Rvec_a|}.
\end{equation}
For the closed-shell implementation in this repository,
\begin{equation}
\rho(\rvec)=2\sum_{n=1}^{N_{\mathrm{occ}}}
|\psi_n(\rvec)|^2.
\end{equation}
The factor two follows from the one-particle density.  Each spatial orbital
has an up-spin and a down-spin copy, both with the same spatial probability:
\begin{align}
\rho(\rvec)
&=\sum_{n,\sigma}|\psi_n(\rvec)\chi_\sigma|^2\nonumber\\
&=\sum_n|\psi_n(\rvec)|^2
\underbrace{\sum_{\sigma}\|\chi_\sigma\|^2}_{2}.
\end{align}

The nuclear potential diverges as \(-Z_a/r\).  Core orbitals therefore vary
rapidly.  Valence orbitals also acquire rapid oscillations near a nucleus
because they must be orthogonal to the occupied core orbitals.  Both effects
require large reciprocal vectors.

\subsection{Why this is especially expensive in plane waves}

At the \(\Gamma\) point, PyPWDFT expands an orbital as
\begin{equation}
\psi_n(\rvec)=\frac{1}{\sqrt{\Omega}}
\sum_{\gvec\in\mathcal B}c_{\gvec n}\e^{\ii\gvec\cdot\rvec},
\qquad
\sum_{\gvec}|c_{\gvec n}|^2=1,
\label{eq:pw-expansion}
\end{equation}
where \(\Omega\) is the cell volume.  Its spherical basis is
\begin{equation}
\mathcal B=
\left\{\gvec:\frac{|\gvec|^2}{2}\le E_{\mathrm{cut}}\right\}.
\label{eq:cutoff}
\end{equation}
The normalization in \cref{eq:pw-expansion} is not an assumption about the
coefficients.  Periodic plane waves satisfy
\begin{equation}
\frac1\Omega\int_\Omega
\e^{\ii(\gvec'-\gvec)\cdot\rvec}\dd^3r
=\delta_{\gvec,\gvec'},
\end{equation}
so substituting the expansion into
\(\int_\Omega|\psi_n|^2\dd^3r=1\) leaves
\(\sum_\gvec|c_{\gvec n}|^2=1\).

Fine real-space features need large \(|\gvec|\), hence a large cutoff and many
plane waves.  The number of basis functions grows approximately as
\begin{equation}
N_{\mathrm{PW}}\sim
\frac{\Omega}{6\pi^2}(2E_{\mathrm{cut}})^{3/2}.
\end{equation}
To see this estimate, one reciprocal-lattice point occupies volume
\((2\pi)^3/\Omega\) in \(\gvec\)-space.  The cutoff keeps a sphere of radius
\(G_{\max}=\sqrt{2E_{\mathrm{cut}}}\); hence
\begin{equation}
N_{\mathrm{PW}}\approx
\frac{(4\pi/3)G_{\max}^3}{(2\pi)^3/\Omega}
=\frac{\Omega G_{\max}^3}{6\pi^2}.
\end{equation}
Surface discreteness makes this an asymptotic count rather than an exact one.
Removing core oscillations can therefore save much more than the cost of
removing a few electrons: it lowers the basis size and every repeated
Hamiltonian application.

\section{From an atom to a pseudopotential}

\subsection{Frozen core}

Choose a division
\begin{equation}
Z= N_{\mathrm{core}}+\zion.
\end{equation}
Only \(\zion\) valence electrons per neutral atom are treated explicitly.
The nucleus plus \(N_{\mathrm{core}}\) frozen electrons becomes an ion of
charge \(+\zion\).  For carbon's bundled \code{C-q4} file,
\(Z=6\), \(N_{\mathrm{core}}=2\), and \(\zion=4\).

``Frozen'' means that core relaxation and explicit core--valence correlation
are not solved during the calculation.  Their mean effect is folded into the
pseudopotential parameters.

\subsection{Pseudo-wavefunctions and matching}

For an isolated reference atom, write a radial orbital as
\(\psi_{nlm}(\rvec)=R_{nl}(r)Y_{lm}(\hat{\rvec})\).  A
norm-conserving construction replaces the rapidly varying all-electron radial
function by a smooth pseudo radial function inside a chosen core radius
\(r_c^l\), while imposing conditions such as
\begin{align}
\widetilde R_{nl}(r)&=R_{nl}^{\mathrm{AE}}(r),
&&r\ge r_c^l,\\
\widetilde\epsilon_{nl}&=\epsilon_{nl}^{\mathrm{AE}},\\
\int_0^{r_c^l}|\widetilde R_{nl}(r)|^2r^2\dd r
&=
\int_0^{r_c^l}|R_{nl}^{\mathrm{AE}}(r)|^2r^2\dd r.
\label{eq:norm-conservation}
\end{align}
The last condition is \emph{norm conservation}.  Together with eigenvalue and
tail matching, it greatly improves the reproduction of scattering over an
energy range near the reference state.  In practical GTH generation,
parameters are optimized against atomic eigenvalues, charge distributions,
and related reference information.  PyPWDFT does not generate these
parameters; it reads an already generated potential.

\paragraph{Why does the norm help scattering?}
Let \(u_E(r)=rR_{El}(r)\) solve the radial equation and let a dot denote
\(\partial/\partial E\).  Subtract the energy derivative of that equation
from the original equation multiplied by \(\dot u_E\).  One integration by
parts gives
\begin{equation}
\int_0^{r_c}u_E(r)^2\dd r
=-\frac12u_E(r_c)^2
\frac{\partial}{\partial E}
\left[\frac{u_E'(r_c)}{u_E(r_c)}\right].
\label{eq:norm-logderivative}
\end{equation}
The ratio \(u_E'/u_E\) at \(r_c\) is the logarithmic derivative that connects
the interior solution to the exterior scattering solution.  If the
all-electron and pseudo functions have the same value, slope, eigenvalue, and
interior norm, \cref{eq:norm-logderivative} also makes their logarithmic
derivatives change identically to first order in energy.  That is the small
mathematical bridge from norm conservation to improved transferability.

\subsection{Transferability}

A potential generated from one atomic configuration is later used in
molecules and solids.  This is the transferability assumption.  It works when
the frozen core remains chemically inert and the chosen reference data cover
the relevant valence physics.  It can fail for unusual oxidation states,
extreme pressure, core spectroscopy, or properties strongly dependent on
near-nuclear density.  Selecting more explicit valence electrons---for
example a \code{q13} rather than \code{q3} potential for Ga---usually improves
the accessible physics but increases cost and often requires a higher cutoff.

\section{Why a non-local operator is needed}

Valence \(s,p,d,\ldots\) states penetrate the core differently.  One purely
local radial function \(V(r)\) generally cannot reproduce all their scattering
properties at once.  A conceptual semi-local form is
\begin{equation}
\widehat V_{\mathrm{SL}}
=\sum_{l,m}|Y_{lm}\rangle V_l(r)\langle Y_{lm}|,
\end{equation}
which explicitly lets each angular momentum see a different radial
interaction.  Such an operator is \emph{non-local}: its action at \(\rvec\)
depends on angular projections of the wavefunction, not only on the value
\(\psi(\rvec)\).

Separable forms rewrite this dependence using a small number of localized
projectors.  For ion \(a\), GTH uses
\begin{equation}
\vnl^{(a)}
=\sum_l\sum_{m=-l}^{l}\sum_{i,j=1}^{n_l}
|p^{a}_{ilm}\rangle h^l_{ij}
\langle p^{a}_{jlm}|.
\label{eq:vnl}
\end{equation}
The \(i,j\) indices enumerate a few radial projectors for the same \(l\);
\(h^l\) is a small real symmetric coupling matrix.  Acting on an orbital,
\begin{equation}
\vnl^{(a)}|\psi\rangle
=\sum_{lmij}|p^{a}_{ilm}\rangle h^l_{ij}
\underbrace{\langle p^{a}_{jlm}|\psi\rangle}_{\text{one scalar overlap}}.
\label{eq:vnl-action}
\end{equation}
This low-rank structure is the central computational advantage.

\section{The GTH analytic form}

GTH means Goedecker--Teter--Hutter.  ``Dual-space Gaussian'' means that the
local smoothing functions and non-local projectors are Gaussian times
polynomials.  They are localized and smooth in real space, while their
Fourier transforms are also analytic Gaussian-polynomial expressions.  That
is an unusually good match to a plane-wave code.

\subsection{Local part in real space}

For one species at the origin, the GTH local potential is
\begin{align}
\vloc(r)
={}&-\frac{\zion}{r}
\erf\left(\frac{r}{\sqrt{2}\,r_{\mathrm{loc}}}\right)
\nonumber\\
&+\e^{-x^2/2}
\left(C_1+C_2x^2+C_3x^4+C_4x^6\right),
\nonumber\\[-2pt]
&\hspace{1em}x=\frac{r}{r_{\mathrm{loc}}}.
\label{eq:vloc-real}
\end{align}
Far from the ion, \(\erf(\cdot)\to1\), the Gaussian term vanishes, and
\(\vloc(r)\to-\zion/r\), as required.  At the origin,
\(\erf(r/\sqrt{2}r_{\mathrm{loc}})/r\) has a finite limit, so the Coulomb
singularity has disappeared.

That finite limit follows directly from
\(\erf z=2(z-z^3/3+\cdots)/\sqrt\pi\):
\begin{equation}
-\frac{\zion}{r}
\erf\left(\frac{r}{\sqrt2r_{\mathrm{loc}}}\right)
=-\zion\sqrt{\frac2\pi}\frac1{r_{\mathrm{loc}}}+O(r^2).
\end{equation}
Therefore
\begin{equation}
V_{\mathrm{loc}}(0)=
-\zion\sqrt{\frac2\pi}\frac1{r_{\mathrm{loc}}}+C_1,
\end{equation}
which makes the contrast with the all-electron \(-Z/r\) divergence explicit.

The five local parameters are \(r_{\mathrm{loc}}\) and up to four \(C_i\).
Unused coefficients are zero-filled by the parser.

\subsection{Local part in reciprocal space}

Use the continuous transform convention
\begin{align}
\widetilde f(\gvec)
&=\int_{\mathbb R^3}f(\rvec)\e^{-\ii\gvec\cdot\rvec}\dd^3r,\\
f(\rvec)
&=\frac{1}{\Omega}\sum_{\gvec}
\widetilde f(\gvec)\e^{\ii\gvec\cdot\rvec}.
\label{eq:fourier-convention}
\end{align}
With \(u=r_{\mathrm{loc}}^2G^2\), \(G=|\gvec|\), the analytic transform of
\cref{eq:vloc-real} for \(G\ne0\) is
\begin{align}
\widetilde V_{\mathrm{loc}}(G)
={}&-\frac{4\pi\zion}{G^2}\e^{-u/2}
\nonumber\\
&+(2\pi)^{3/2}r_{\mathrm{loc}}^3\e^{-u/2}
\Big[
C_1+C_2(3-u)
\nonumber\\
&\quad+C_3(15-10u+u^2)
\nonumber\\
&\quad+C_4(105-105u+21u^2-u^3)
\Big].
\label{eq:vloc-fourier}
\end{align}

\paragraph{Deriving the Gaussian-polynomial terms.}
Write \(\sigma=r_{\mathrm{loc}}\).  The three-dimensional Gaussian transform
factorizes into three elementary one-dimensional integrals:
\begin{align}
F_0(G)
&=\prod_{\alpha=x,y,z}
\int_{-\infty}^{\infty}
\e^{-r_\alpha^2/(2\sigma^2)}
\e^{-\ii G_\alpha r_\alpha}\dd r_\alpha\nonumber\\
&=(2\pi)^{3/2}\sigma^3\e^{-\sigma^2G^2/2}.
\end{align}
Since
\(-\nabla_{\gvec}^2\e^{-\ii\gvec\cdot\rvec}
=r^2\e^{-\ii\gvec\cdot\rvec}\), multiplication by \(r^2\) before
transforming is the same as applying \(-\nabla_\gvec^2\) afterwards.  Thus
\begin{align}
\mathcal F\{x^2\e^{-x^2/2}\}
&=(3-u)F_0,\nonumber\\
\mathcal F\{x^4\e^{-x^2/2}\}
&=(15-10u+u^2)F_0,\nonumber\\
\mathcal F\{x^6\e^{-x^2/2}\}
&=(105-105u+21u^2-u^3)F_0,
\end{align}
where the powers of \(\sigma\) cancel because \(x=r/\sigma\).

For the smoothed Coulomb term, regard
\(-\zion\erf(r/\sqrt2\sigma)/r\) as the potential of a normalized Gaussian
ionic charge.  Applying Poisson's equation and then Fourier transforming gives
\begin{align}
-G^2\widetilde V_{\mathrm C}(G)
&=4\pi\zion\e^{-\sigma^2G^2/2},\nonumber\\
\widetilde V_{\mathrm C}(G)
&=-\frac{4\pi\zion}{G^2}\e^{-u/2}.
\end{align}
Adding this to the three polynomial transforms yields
\cref{eq:vloc-fourier}, term by term.  Those same terms appear directly in
\code{\_build\_local\_potential()}.

\subsection{The special \texorpdfstring{\(G=0\)}{G=0} value}

The periodic Coulomb zero mode is singular.  Expand the first term:
\begin{equation}
-\frac{4\pi\zion}{G^2}\e^{-r_{\mathrm{loc}}^2G^2/2}
=-\frac{4\pi\zion}{G^2}
+2\pi\zion r_{\mathrm{loc}}^2+O(G^2).
\end{equation}
Here we used only
\(\e^{-u/2}=1-u/2+O(u^2)\), with
\(u=r_{\mathrm{loc}}^2G^2\).  The \(1/G^2\) multiplying the linear term in
\(u\) cancels, leaving the finite second term.
The code removes the divergent Coulomb component, which is handled through
the neutral periodic electrostatic convention, but retains the finite limit.
Consequently it assigns
\begin{align}
\widetilde V_{\mathrm{loc}}(0)
={}&2\pi\zion r_{\mathrm{loc}}^2
\nonumber\\
&+(2\pi)^{3/2}r_{\mathrm{loc}}^3
(C_1+3C_2+15C_3+105C_4).
\label{eq:gzero}
\end{align}
Thus ``the \(G=0\) value is zero'' would be an incorrect description of this
implementation.

\subsection{Many ions and the structure factor}

Translation by \(\Rvec_a\) contributes a phase.  In a conventional inverse
transform one writes
\begin{equation}
\widetilde V_{\mathrm{ion}}(\gvec)
=\sum_a\widetilde V_{\mathrm{loc}}^{s_a}(G)
\e^{-\ii\gvec\cdot\Rvec_a},
\label{eq:structure-factor}
\end{equation}
where \(s_a\) is the species.  PyPWDFT groups atoms by species, evaluates a
species structure factor, and transforms once:
\begin{equation}
S_s(\gvec)=\sum_{a:s_a=s}\e^{+\ii\gvec\cdot\Rvec_a}.
\end{equation}
The phase is just the Fourier translation theorem.  With
\(\rvec'=\rvec-\Rvec_a\),
\begin{align}
\int V(\rvec-\Rvec_a)\e^{-\ii\gvec\cdot\rvec}\dd^3r
&=\int V(\rvec')\e^{-\ii\gvec\cdot(\rvec'+\Rvec_a)}
\dd^3r'\nonumber\\
&=\e^{-\ii\gvec\cdot\Rvec_a}\widetilde V(\gvec).
\end{align}
Why the plus sign?  NumPy's \code{fftn} uses the negative exponential.  The
combination of a plus-sign structure factor and \code{fftn} produces
\(\vloc(\rvec-\Rvec_a)\).  Re-indexing \(\gvec\to-\gvec\) shows its
equivalence to \cref{eq:structure-factor}.  For a real potential, the
reciprocal coefficients obey Hermitian symmetry, and the tiny residual
imaginary part is discarded.

\section{The GTH projectors}

\subsection{Real-space definition}

A normalized GTH radial projector can be written
\begin{equation}
p_i^l(r)=
\frac{\sqrt{2}\,r^{l+2(i-1)}
\e^{-r^2/(2r_l^2)}}
{r_l^{\,l+(4i-1)/2}
\sqrt{\Gamma\!\left(l+\frac{4i-1}{2}\right)}}.
\label{eq:radial-projector}
\end{equation}
Together with a real spherical harmonic,
\begin{equation}
p_{ilm}(\rvec)=p_i^l(r)Y_{lm}^{\mathrm R}(\hat{\rvec}),
\end{equation}
it is smooth and localized near the ion.  A projector on atom \(a\) is the
translated function \(p^a_{ilm}(\rvec)=p_{ilm}(\rvec-\Rvec_a)\).

\paragraph{Checking the normalization.}
Squaring \cref{eq:radial-projector} leaves the radial integral
\begin{align}
I&=\int_0^\infty r^{2l+4i-2}\e^{-r^2/r_l^2}\dd r\nonumber\\
&=\frac12r_l^{2l+4i-1}
\Gamma\left(l+\frac{4i-1}{2}\right),
\end{align}
where \(t=r^2/r_l^2\) turns the integral into the definition of the gamma
function.  The squared prefactor in \cref{eq:radial-projector} is exactly
\(I^{-1}\), so
\(\int_0^\infty|p_i^l(r)|^2r^2\dd r=1\).

\subsection{Plane-wave overlaps}

The code never samples \cref{eq:radial-projector} in real space.  It directly
constructs the plane-wave overlap
\begin{align}
\beta^{a,lm}_{\gvec i}
&=\langle\gvec|p^a_{ilm}\rangle\\
&=(-\ii)^lY_{lm}^{\mathrm R}(\hat{\gvec})\,
P_i^l(G)\,\e^{-\ii\gvec\cdot\Rvec_a},
\label{eq:beta}
\end{align}
on only the active plane waves from \cref{eq:cutoff}.  The factor
\(\Omega^{-1/2}\) needed for normalized plane waves is contained in the
radial factor below.

The three factors in \cref{eq:beta} follow separately.  First, the translated
projector supplies \(\e^{-\ii\gvec\cdot\Rvec_a}\) by the translation theorem.
Second, the plane-wave expansion
\begin{equation}
\e^{-\ii\gvec\cdot\rvec}
=4\pi\sum_{l,m}(-\ii)^l
j_l(Gr)Y_{lm}(\hat{\gvec})Y_{lm}^*(\hat{\rvec})
\end{equation}
and angular orthogonality select only the matching \(l,m\).  Finally the
remaining radial integral is
\begin{equation}
P_i^l(G)\propto
\frac{4\pi}{\sqrt\Omega}
\int_0^\infty r^2p_i^l(r)j_l(Gr)\dd r.
\end{equation}
A Gaussian times \(r^{l+2(i-1)}\) transforms into a Gaussian times
\(G^l\) and a polynomial in \(G^2\), producing the forms below.

Let \(q=Gr_l\) and
\begin{equation}
A_l(G)=4\pi^{5/4}
\sqrt{\frac{2^{l+1}r_l^{2l+3}}{\Omega}}\,
\e^{-q^2/2}.
\end{equation}
The implemented radial transforms are
\begin{align}
P_1^0&=A_0,&
P_2^0&=\frac{2(3-q^2)}{\sqrt{15}}A_0,\nonumber\\
P_3^0&=\frac{4(15-10q^2+q^4)}
{3\sqrt{105}}A_0,\nonumber\\
P_1^1&=\frac{G}{\sqrt3}A_1,&
P_2^1&=\frac{2G(5-q^2)}{\sqrt{105}}A_1,\nonumber\\
P_3^1&=\frac{4G(35-14q^2+q^4)}
{3\sqrt{1155}}A_1,\nonumber\\
P_1^2&=\frac{G^2}{\sqrt{15}}A_2,&
P_2^2&=\frac{2G^2(7-q^2)}
{3\sqrt{105}}A_2,\nonumber\\
P_1^3&=\frac{G^3}{\sqrt{105}}A_3.
\label{eq:implemented-projectors}
\end{align}
These branches are literal translations of
\code{\_evaluate\_projector()}.  The parser permits at most three projectors
per \(l\), but the evaluator implements only the combinations in
\cref{eq:implemented-projectors}.  The bundled files are compatible with
that set; an external file requesting, for example, a third \(d\) projector
would parse and then raise an error while projectors are built.

\subsection{Real spherical harmonics}

For every \(l\), the implementation loops over all
\(m=-l,\ldots,l\).  It evaluates a real, orthonormal tesseral basis.  For
example,
\begin{align}
Y_{00}^{\mathrm R}&=\frac{1}{2}\sqrt{\frac1\pi},\\
Y_{10}^{\mathrm R}&=\frac{1}{2}\sqrt{\frac3\pi}\cos\theta,\\
Y_{11}^{\mathrm R}&=\frac{1}{2}\sqrt{\frac3\pi}
\sin\theta\cos\phi,\\
Y_{1,-1}^{\mathrm R}&=\frac{1}{2}\sqrt{\frac3\pi}
\sin\theta\sin\phi.
\end{align}
Changing between real and complex spherical harmonics is only a unitary
change of basis within a complete \(l\) channel.  Summing all \(m\) values
preserves rotational invariance.

At \(G=0\), the direction \(\hat{\gvec}\) is undefined.  The code assigns a
harmless angular direction for \(l>0\); the explicit factors
\(G^l\) make those projector values zero.  The \(s\) harmonic is constant.

\section{The non-local operation as linear algebra}

For a fixed atom, \(l\), and \(m\), collect the projector columns into
\begin{equation}
\begin{aligned}
B&=
\begin{bmatrix}
\vert&&\vert\\
\bm\beta_1&\cdots&\bm\beta_{n_l}\\
\vert&&\vert
\end{bmatrix}
&&\in\mathbb C^{N_{\mathrm{PW}}\times n_l},\\
H&=h^l
&&\in\mathbb R^{n_l\times n_l}.
\end{aligned}
\end{equation}
For \(N_s\) trial orbitals, collect their coefficients in
\(C\in\mathbb C^{N_{\mathrm{PW}}\times N_s}\).  One channel is applied as
\begin{equation}
\boxed{\quad C\longmapsto B\big[H(B^\dagger C)\big].\quad}
\label{eq:bhb}
\end{equation}
This matrix expression is obtained by inserting plane-wave resolutions of the
identity on both sides of \cref{eq:vnl}.  In indices, one channel gives
\begin{align}
(\vnl C)_{\gvec n}
&=\sum_{\gvec'}\sum_{i,j}
\langle\gvec|p_i\rangle h_{ij}
\langle p_j|\gvec'\rangle C_{\gvec'n}\nonumber\\
&=\sum_i B_{\gvec i}\sum_jH_{ij}
\sum_{\gvec'}B_{\gvec'j}^*C_{\gvec'n}.
\end{align}
The innermost sum is \((B^\dagger C)_{jn}\); regrouping the other two sums
then gives \(B[H(B^\dagger C)]\).

This is exactly the NumPy expression
\begin{lstlisting}[language=Python]
result += beta @ (
    coupling @ (beta.conj().T @ coefficients)
)
\end{lstlisting}
Read it from right to left:
\begin{enumerate}
  \item \(B^\dagger C\): overlap each state with each localized projector;
  \item \(H(B^\dagger C)\): mix radial projectors within the \(l\) channel;
  \item \(B[\cdots]\): scatter the small result back over all plane waves.
\end{enumerate}

No dense \(N_{\mathrm{PW}}\times N_{\mathrm{PW}}\) matrix is formed.  If
\(n_l\) is small, the leading work is proportional to
\(N_{\mathrm{PW}}n_lN_s\).  The CPU and GPU paths use the same algebra.  The
GPU path caches device copies of \(B\) and \(H\) so iterative eigensolver
calls do not repeatedly transfer them.

\subsection{Why the operator is Hermitian}

The parser reads the upper triangle of every coupling matrix and reflects it:
\(H=H^\mathsf T\).  Therefore
\begin{equation}
(BHB^\dagger)^\dagger
=BH^\dagger B^\dagger=BHB^\dagger.
\end{equation}
Hermiticity guarantees real expectation values and is required by the
Hermitian eigensolver.  The test
\code{test\_nonlocal\_operator\_is\_hermitian} checks numerically that
\begin{equation}
\langle x|\vnl y\rangle
=\langle\vnl x|y\rangle.
\end{equation}

\section{Decoding a parameter file}

Consider \path{pypwdft/psp/pade/C-q4}:
\begin{lstlisting}
C GTH-PADE-q4 GTH-LDA-q4
    2    2
    0.34883045  2  -8.51377110  1.22843203
    2
    0.30455321  1   9.52284179
    0.23267730  0
\end{lstlisting}
It means:
\begin{enumerate}
  \item The header identifies carbon, the GTH-PADE/LDA family, and the
  \(q4\) valence partition.
  \item The occupation integers sum to \(\zion=2+2=4\).  The parser needs
  only their sum.
  \item \(r_{\mathrm{loc}}=0.34883045\) bohr.  Two local coefficients follow:
  \(C_1=-8.51377110\), \(C_2=1.22843203\); \(C_3=C_4=0\).
  \item The next integer is \(l_{\max}=2\) in the file convention used by
  the parser: there are two channel records, \(l=0\) and \(l=1\).
  It is a count, not an inclusive highest-\(l\) index.
  \item The \(s\) channel has \(r_0=0.30455321\), one projector, and
  \(h_{11}^{0}=9.52284179\).
  \item The \(p\) channel records \(r_1=0.23267730\) but zero projectors, so
  it contributes nothing to \(\vnl\).
\end{enumerate}
Thus one carbon atom using this file has one actual non-local projector
column.  Hydrogen's bundled PADE file has \(l_{\max}=0\), so its
non-local contribution is identically zero; this is why the H\(_2\)
integration test reports \(E_{\mathrm{nl}}=0\).

For a channel with several projectors, the file stores the upper triangle of
\(h^l\) over several lines.  The parser inserts those entries and mirrors
them.  Arrays are allocated for \(l=0,1,2,3\), at most three radial
projectors per channel, and four local coefficients.

\section{Repository Fourier and grid conventions}

\subsection{The reciprocal grid}

For a cubic cell of edge \(L\), reciprocal vectors have components
\begin{equation}
G_\alpha=\frac{2\pi n_\alpha}{L},
\qquad n_\alpha\in\mathbb Z.
\end{equation}
\code{PeriodicSystem.\_\_build\_fft\_vectors()} obtains the FFT-ordered
integers from \code{numpy.fft.fftfreq}.  The orbital mask is precisely
\(\tfrac12G^2\le E_{\mathrm{cut}}\).  The working FFT grid is large enough
for a density cutoff of at least \(4E_{\mathrm{cut}}\), because multiplying
two orbitals can double the largest reciprocal wavevector and therefore
quadruple its kinetic-energy scale.

Explicitly, a density product contains
\begin{equation}
\psi_n^*(\rvec)\psi_n(\rvec)
=\frac1\Omega\sum_{\gvec,\gvec'}
c_{\gvec n}^*c_{\gvec'n}
\e^{\ii(\gvec'-\gvec)\cdot\rvec}.
\end{equation}
If both orbital vectors are bounded by \(G_{\max}\), their difference can
reach \(2G_{\max}\).  Its kinetic-energy scale is
\(\tfrac12(2G_{\max})^2=4E_{\mathrm{cut}}\), which derives the factor four.

\subsection{Normalization in the discrete FFT}

NumPy defines an unnormalized forward transform and a \(1/N^3\)-normalized
inverse transform.  If an active coefficient vector \(c_{\gvec}\) is embedded
in the full \(N^3\) grid, PyPWDFT computes
\begin{equation}
\psi(\rvec_j)=
\frac{\operatorname{ifftn}(c)_j}{C_t},
\qquad
C_t=\frac{\sqrt{\Omega}}{N^3}.
\end{equation}
Indeed, NumPy's definition is
\begin{equation}
\operatorname{ifftn}(c)_j
=\frac1{N^3}\sum_\gvec c_\gvec
\e^{\ii\gvec\cdot\rvec_j}.
\end{equation}
Dividing by \(C_t=\sqrt\Omega/N^3\) cancels \(N^{-3}\) and leaves
\(\Omega^{-1/2}\sum_\gvec c_\gvec\e^{\ii\gvec\cdot\rvec_j}\).
This recovers exactly the normalized plane-wave expression
\cref{eq:pw-expansion}.  The coefficient-space eigensolver can consequently
use the ordinary Euclidean inner product.

\section{How the pseudopotential enters the Hamiltonian}

With a GTH ion model, the Kohn--Sham Hamiltonian used by the solver is
\begin{equation}
\widehat H[\rho]
=-\frac12\nabla^2
+V_{\mathrm{loc}}^{\mathrm{ions}}(\rvec)
+V_{\mathrm H}[\rho](\rvec)
+V_{\mathrm{xc}}[\rho](\rvec)
+\vnl.
\label{eq:ks-psp}
\end{equation}
It contains three computational patterns:
\begin{enumerate}
  \item \textbf{Kinetic:} diagonal in reciprocal space,
  \((Tc)_{\gvec}=\tfrac12G^2c_{\gvec}\).
  \item \textbf{Local potentials:} transform \(c\) to real space, multiply
  pointwise by \(V_{\mathrm{loc}}^{\mathrm{ions}}+V_{\mathrm H}
  +V_{\mathrm{xc}}\), and transform back.
  \item \textbf{Non-local potential:} apply the projector products in
  \cref{eq:bhb} directly in the active reciprocal basis.
\end{enumerate}
The kinetic result is immediate because
\begin{equation}
-\frac12\nabla^2
\left(\frac{\e^{\ii\gvec\cdot\rvec}}{\sqrt\Omega}\right)
=\frac{G^2}{2}
\left(\frac{\e^{\ii\gvec\cdot\rvec}}{\sqrt\Omega}\right).
\end{equation}
Every plane wave is therefore an eigenfunction of the kinetic operator; only
its own coefficient is multiplied, so no FFT or dense matrix is needed.

In condensed pseudocode, the matrix-free operation in \code{LinOpH} is
\begin{lstlisting}[language=Python]
psi_G[mask] = c
kinetic = 0.5 * G2[mask] * c
local = fftn(ifftn(psi_G) * Veff)[mask]
result = kinetic + local + apply_nonlocal(c)
\end{lstlisting}
The iterative \code{eigsh} solver requests only this action; it never asks
PyPWDFT to assemble the full Hamiltonian matrix.

\section{Electron count and energy bookkeeping}

\subsection{Which electron count is used?}

Without a pseudopotential, the low-level system sums nuclear atomic numbers.
With GTH, \code{GTHPseudopotential.nelec} instead sums the \(\zion\) values
from the selected files:
\begin{equation}
N_e=\sum_a Z_{\mathrm{ion},a}.
\end{equation}
This describes a neutral cell.  The present solver then requires \(N_e\) to
be even and occupies \(N_e/2\) spatial orbitals twice.  There is no public
charged-cell, spin-polarized, fractional-occupation, or smearing path in this
implementation.

\subsection{Local electron--ion energy}

For the converged density,
\begin{equation}
E_{\mathrm{e-i}}^{\mathrm{loc}}
=\int_\Omega \rho(\rvec)
V_{\mathrm{loc}}^{\mathrm{ions}}(\rvec)\dd^3r.
\end{equation}
Partitioning a uniform cell into \(N^3\) equal grid boxes turns the integral
into the Riemann sum
\begin{equation}
E_{\mathrm{e-i}}^{\mathrm{loc}}
\approx\sum_j\rho_jV_j\Delta V,
\qquad
\Delta V=\frac{\Omega}{N^3},
\end{equation}
implemented by an \code{einsum}.  The equality becomes exact for Fourier
components represented without aliasing on the grid.

\subsection{Non-local energy}

For doubly occupied normalized orbitals,
\begin{align}
E_{\mathrm{e-i}}^{\mathrm{nl}}
&=2\sum_{n=1}^{N_{\mathrm{occ}}}
\langle\psi_n|\vnl|\psi_n\rangle\\
&=2\,\mathrm{Re}\!
\sum_{n,\gvec}c_{\gvec n}^*
(\vnl C)_{\gvec n}.
\label{eq:enl}
\end{align}
The second line is simply the coefficient-space inner product:
\begin{equation}
\langle\psi_n|\vnl|\psi_n\rangle
=\sum_\gvec c_{\gvec n}^*(\vnl\bm c_n)_\gvec.
\end{equation}
Summing the two spin states of each occupied spatial orbital supplies the
factor two.  The explicit real part removes only roundoff-level imaginary
noise; Hermiticity makes the exact expectation value real.
The factor two is spin occupancy, not part of the pseudopotential itself.
\code{nonlocal\_energy()} and the SCF code implement the same expression.

\subsection{Ion--ion energy}

The residual ions carry charges \(+\zion\), not bare nuclear charges \(+Z\).
The Ewald derivation starts by adding and subtracting a Gaussian screening
charge, equivalently splitting
\begin{equation}
\frac1r=
\underbrace{\frac{\erfc(\sqrt\alpha r)}{r}}_{\text{short-ranged}}
+\underbrace{\frac{\erf(\sqrt\alpha r)}{r}}_{\text{smooth, long-ranged}}.
\end{equation}
The first term converges rapidly in direct-lattice vectors \(\mathbf T\).
The second has the reciprocal transform
\(4\pi\e^{-G^2/(4\alpha)}/G^2\), so it converges rapidly in \(\gvec\)-space.
One must then subtract each Gaussian's self-interaction and specify the
periodic zero-mode/background convention.  These four pieces give
\begin{align}
E_{\mathrm{ion-ion}}
={}&\frac12\sum_{a,b,\mathbf T}'\!
\frac{Z_a^{\mathrm{ion}}Z_b^{\mathrm{ion}}\,
\erfc(\sqrt{\alpha}|\Rvec_{ab}+\mathbf T|)}
{|\Rvec_{ab}+\mathbf T|}
\nonumber\\
&+\frac{2\pi}{\Omega}\sum_{\gvec\ne0}
\frac{\e^{-G^2/(4\alpha)}}{G^2}
\left|\sum_aZ_a^{\mathrm{ion}}
\e^{\ii\gvec\cdot\Rvec_a}\right|^2
\nonumber\\
&-\sqrt{\frac{\alpha}{\pi}}\sum_a
(Z_a^{\mathrm{ion}})^2
-\frac{\pi}{2\alpha\Omega}
\left(\sum_aZ_a^{\mathrm{ion}}\right)^2.
\label{eq:ewald}
\end{align}
The prime excludes the \(a=b,\mathbf T=0\) self pair.  The last term is the
uniform-background/electroneutrality convention used by
\code{calculate\_ewald\_sum()}.  The GTH wrapper calls that method with
\code{valence\_charges}; the all-electron path instead uses nuclear charges.

\subsection{Total energy}

The reported energy is
\begin{align}
E_{\mathrm{tot}}
={}&E_{\mathrm{kin}}
+E_{\mathrm{e-i}}^{\mathrm{loc}}
+E_{\mathrm{e-i}}^{\mathrm{nl}}
\nonumber\\
&+E_{\mathrm H}+E_{\mathrm{xc}}
+E_{\mathrm{ion-ion}}.
\label{eq:etot}
\end{align}
The non-local energy must be added explicitly because it is not represented
on the real-space potential grid.  It would be an error either to omit it or
to count it inside the local electron--ion integral.

\section{End-to-end execution trace}

The complete life of a pseudopotential in one PyPWDFT calculation is:
\begin{enumerate}
  \item The high-level \code{PWDFT} object converts element symbols to nuclear
  atomic numbers in a \code{PeriodicSystem}.
  \item \code{\_build\_pseudopotential()} chooses GTH-PADE for LDA or GTH-PBE
  for PBE.  An explicit \code{GTH(charges=\{...\})} can request another
  bundled \(q\) value.
  \item \code{GTHPseudopotential} converts atomic numbers back to element
  symbols, finds one file per distinct species, and parses its parameters.
  By default, the available file with the smallest \(q\) is selected.
  \item It assigns one \(\zion\) to every atom and sums them to obtain the
  explicit electron count.
  \item \code{rebuild()} caches the full real-space local ionic field and all
  active-basis projector matrices.  Both depend on atomic positions.
  \item Before the SCF loop, the solver obtains the local grid, a callable
  non-local operator, and the valence-charge Ewald energy.
  \item During each iterative eigenvalue solve, \code{LinOpH} adds kinetic,
  local, and non-local actions according to \cref{eq:ks-psp}.
  \item Occupied orbitals create a new density.  Hartree and XC quantities
  change with that density; the ionic local and non-local terms do not.
  \item Energy evaluation uses \cref{eq:enl} and \cref{eq:etot}.
  \item After SCF convergence, the code performs a final solve and recomputes
  every energy term from a mutually consistent final set of orbitals and
  density.
\end{enumerate}

\section{A map from mathematics to source}

\begin{table}[ht]
\centering
\footnotesize
\begin{tabular}{@{}p{0.28\columnwidth}p{0.66\columnwidth}@{}}
\toprule
Concept & Current implementation location\\
\midrule
Plane-wave grid and cutoff &
\path{pypwdft/psystem.py}: \code{\_\_build\_fft\_vectors},
\code{\_\_grid\_size}, and \code{\_\_pw\_mask}\\
Parameter selection &
\path{pypwdft/api.py}: \code{GTH.\_build} and
\code{\_build\_pseudopotential}\\
Parameter parsing &
\path{pypwdft/gth.py}: \code{\_find\_file} and \code{\_read\_gth}\\
Local equations
\cref{eq:vloc-fourier,eq:gzero} &
\path{GTHPseudopotential._build_local_potential}\\
Projector equations
\cref{eq:beta,eq:implemented-projectors} &
\code{\_build\_projectors}, \code{\_evaluate\_projector}, and
\code{\_real\_spherical\_harmonic}\\
Separable action
\cref{eq:bhb} &
\code{apply\_nonlocal} and \code{nonlocal\_operator}\\
Hamiltonian
\cref{eq:ks-psp} &
\path{pypwdft/pypwdft.py}: \code{LinOpH.\_matvec} and GPU
\code{\_CuPyHamiltonian.matvec}\\
Energy
\cref{eq:enl,eq:etot} &
\code{PyPWDFT.scf}, \code{nonlocal\_energy}, and \code{ionic\_energy}\\
Executable checks &
\path{tests/test_gth.py}\\
\bottomrule
\end{tabular}
\caption{A reading guide for the repository.}
\label{tab:code-map}
\end{table}

\section{What the tests prove---and what they do not}

\path{tests/test_gth.py} provides useful executable statements of the
mathematics:
\begin{itemize}
  \item Parsing carbon yields \(\zion=4\), the documented
  \(r_{\mathrm{loc}}\), and local coefficients.
  \item The bundled PBE family gives a different, expected
  \(r_{\mathrm{loc}}\).
  \item Random-vector inner products verify Hermiticity of \(\vnl\).
  \item \code{nonlocal\_energy()} matches an independently written
  expectation value.
  \item A Hamiltonian with zero local potential equals kinetic plus non-local
  action.
  \item Moving atoms invalidates cached phases and potentials until
  \code{rebuild()} is called.
  \item The Ewald routine produces a different answer when passed valence
  rather than nuclear charges.
  \item Small H\(_2\) SCF calculations exercise both LDA/PADE and PBE/PBE
  integration.
\end{itemize}

These tests establish internal consistency.  They do not by themselves
establish transferability or reproduce the atomic scattering tests used to
generate a potential.  Scientific validation additionally requires
comparison with trusted calculations or experiment and convergence with
respect to cutoff, cell size, and other numerical settings.

\section{Choosing and converging a potential}

\subsection{Match the XC family}

The bundled \code{pade} parameters are intended for LDA and the \code{pbe}
parameters for PBE.  The high-level API enforces that match for bundled data.
The name ``PADE'' here labels the GTH LDA parameter family; it is not the
plane-wave basis and not an alternative SCF algorithm.

\subsection{Understand \texorpdfstring{\(qN\)}{qN}}

The suffix \(qN\) is the number of explicit valence electrons per neutral
atom.  It is not the atom's net oxidation state and not its nuclear atomic
number.  The repository selects the lowest available \(q\) by default.
\begin{lstlisting}[language=Python]
GTH(charges={"Ga": 13})
\end{lstlisting}
requests \code{Ga-q13} rather than the default lower-valence Ga file.
Different \(q\) choices define different Hamiltonians and electron counts;
their raw total energies should not be mixed casually.

\subsection{Converge the plane-wave cutoff}

Gaussian pseudopotentials make modest cutoffs possible, but do not make a
particular cutoff automatically accurate.  Increase \(E_{\mathrm{cut}}\)
until energy differences or other target observables change less than the
desired tolerance.  A harder, smaller-core, higher-\(q\) potential commonly
needs a larger cutoff.  The \(4E_{\mathrm{cut}}\) density-grid rule prevents
aliasing of orbital products; it is distinct from convergence of the orbital
basis itself.

\subsection{Converge the periodic cell for molecules}

PyPWDFT uses periodic boundary conditions and currently only the
\(\Gamma\) point.  A molecule is therefore an infinite array of molecular
images.  Increase the vacuum until image interactions are acceptably small.
Pseudopotentials remove core cost; they do not remove finite-cell
electrostatics.

\section{Scope and limitations of this implementation}

The current code is educational and deliberately compact.  When explaining
it, keep these boundaries visible:
\begin{itemize}
  \item cubic periodic cells and \(\Gamma\)-point orbitals only;
  \item spin-unpolarized, doubly occupied orbitals and an even electron count;
  \item norm-conserving separable GTH form, not ultrasoft pseudopotentials or
  the projector augmented-wave (PAW) method;
  \item no non-linear core correction in the shown GTH energy functional;
  \item no spin--orbit operator; heavy-element parameters may incorporate
  scalar-relativistic fitting, but the Hamiltonian has no spinor channels;
  \item no forces, stress, geometry optimization, or molecular dynamics;
  cached geometry-dependent quantities require an explicit \code{rebuild()};
  \item a fixed set of analytic projector branches, listed in
  \cref{eq:implemented-projectors};
  \item neutral cells as determined by the selected \(\zion\) values.
\end{itemize}

\section{Common misconceptions}

\paragraph{``The core electrons are deleted.''}
They are removed from the explicit variational problem, but their average
effect and the nuclear field are represented by fitted potential parameters.

\paragraph{``A pseudopotential is just \( -Z/r\) with the tip rounded.''}
That describes only the local intuition.  Angular-momentum-dependent
scattering requires the non-local projector operator.

\paragraph{``Non-local means long-ranged.''}
Not here.  The projectors are spatially localized Gaussians.  ``Non-local''
means the operator is an integral/projection operator rather than pointwise
multiplication.

\paragraph{``The pseudo-wavefunction is the true wavefunction everywhere.''}
It is intentionally different inside the core.  Near-nuclear density,
hyperfine quantities, and core spectra generally need reconstruction or an
all-electron/PAW treatment.

\paragraph{``A larger \(q\) is always better.''}
It treats more shells explicitly and may improve transferability for a target
property, but costs more and can be harder.  ``Better'' is property- and
validation-dependent.

\paragraph{``The sum of Kohn--Sham eigenvalues is the total energy.''}
It double-counts some density-dependent interactions and misses the required
DFT corrections.  PyPWDFT assembles the explicit components in
\cref{eq:etot}.

\paragraph{``The local and non-local energies are optional diagnostics.''}
Both belong to the Hamiltonian and total-energy functional.  Splitting them
is a representation choice; only their consistent combination with the
ion--ion and electronic terms is physically meaningful.

\section{Two additional algebraic checks}

\subsection{Local-potential matrix elements}

From normalized plane waves and \cref{eq:fourier-convention},
\begin{align}
\langle\gvec|\vloc|\gvec'\rangle
&=\frac1\Omega\int_\Omega
\e^{-\ii\gvec\cdot\rvec}\vloc(\rvec)
\e^{\ii\gvec'\cdot\rvec}\dd^3r\\
&=\frac1\Omega
\widetilde V_{\mathrm{loc}}(\gvec-\gvec').
\end{align}
The matrix is dense in reciprocal space.  Rather than form it, PyPWDFT
transforms the orbital to real space, multiplies pointwise, and transforms
back.  This FFT convolution is the local counterpart to the low-rank
non-local operation.

\subsection{Non-local energy in projector overlaps}

Insert \cref{eq:vnl} into one expectation value:
\begin{equation}
\langle\psi|\vnl|\psi\rangle
=\sum_{a,l,m}\bm q_{a,l,m}^\dagger
h^l\bm q_{a,l,m},
\qquad
q_i=\langle p^a_{ilm}|\psi\rangle.
\end{equation}
This shows both why the operation is cheap and why off-diagonal \(h_{ij}^l\)
values matter: they couple different radial projectors in the same angular
channel.

\section{How to explain the code aloud}

A concise but accurate explanation is:
\begin{quote}\small
The code uses norm-conserving GTH pseudopotentials to freeze atomic core
electrons.  Each parameter file supplies a valence ionic charge, a smooth
local Gaussian-polynomial potential, and small angular-momentum projector
matrices.  The local potential has an analytic Fourier transform, so the code
evaluates it on the reciprocal FFT grid, includes atomic positions through
structure factors, and transforms it once to real space.  The non-local part
is kept in separable form.  For each atom and \(l,m\) channel, a matrix of
plane-wave projector overlaps \(B\) is built analytically, and its action is
\(B h B^\dagger c\).  The Hamiltonian combines diagonal reciprocal-space
kinetic energy, FFT-based multiplication by all local potentials, and that
low-rank projector action.  The total energy separately includes local and
non-local electron--ion terms and an Ewald energy computed with valence
ionic charges.
\end{quote}

If more detail is requested, explain in this order:
\begin{enumerate}
  \item why all-electron core oscillations are expensive in plane waves;
  \item frozen core, norm conservation, and transferability;
  \item why different \(l\) channels force a non-local term;
  \item the GTH Gaussian form in real and reciprocal space;
  \item the \(B h B^\dagger\) matrix product and its dimensions;
  \item parameter parsing, electron count, and energy bookkeeping;
  \item numerical and physical limitations.
\end{enumerate}

\section{Checklist for reading a new code path}

When encountering a pseudopotential implementation, ask:
\begin{enumerate}
  \item What Fourier and plane-wave normalization conventions are used?
  \item What is the active plane-wave mask and cutoff?
  \item Which electrons are frozen, and where is \(\zion\) obtained?
  \item How is the local \(G=0\) component defined?
  \item What are the projector indices, shapes, phases, and harmonics?
  \item Is the coupling matrix explicitly Hermitian?
  \item Are atomic translations applied with a consistent phase?
  \item Is the non-local operator included in both the Hamiltonian and energy?
  \item Does the ion--ion term use valence rather than nuclear charges?
  \item Does the parameter family match the XC functional?
  \item Which projector combinations and angular momenta are supported?
  \item What changes must invalidate geometry-dependent caches?
\end{enumerate}
Answering those questions for PyPWDFT reproduces essentially the whole
implementation described in this note.

\section*{Notation summary}

\begin{center}
\footnotesize
\begin{tabular}{@{}ll@{}}
\toprule
Symbol & Meaning\\
\midrule
\(\Omega\) & periodic cell volume\\
\(\Rvec_a\) & position of ion \(a\)\\
\(\gvec\) & reciprocal lattice vector\\
\(E_{\mathrm{cut}}\) & orbital plane-wave cutoff\\
\(Z\) & nuclear atomic number\\
\(\zion\) & explicit valence/ionic charge\\
\(r_{\mathrm{loc}},C_i\) & local GTH parameters\\
\(r_l\) & projector radius for channel \(l\)\\
\(i,j\) & radial-projector indices\\
\(l,m\) & angular-momentum indices\\
\(h^l\) & small projector coupling matrix\\
\(B\) or \(\beta\) & plane-wave projector overlaps\\
\bottomrule
\end{tabular}
\end{center}

\begin{thebibliography}{9}

\bibitem{GTH1996}
S.~Goedecker, M.~Teter, and J.~Hutter,
``Separable dual-space Gaussian pseudopotentials,''
\emph{Physical Review B} \textbf{54}, 1703--1710 (1996).
\href{https://doi.org/10.1103/PhysRevB.54.1703}
{doi:10.1103/PhysRevB.54.1703}.

\bibitem{HGH1998}
C.~Hartwigsen, S.~Goedecker, and J.~Hutter,
``Relativistic separable dual-space Gaussian pseudopotentials from H to Rn,''
\emph{Physical Review B} \textbf{58}, 3641--3662 (1998).
\href{https://doi.org/10.1103/PhysRevB.58.3641}
{doi:10.1103/PhysRevB.58.3641}.

\bibitem{Krack2005}
M.~Krack,
``Pseudopotentials for H to Kr optimized for gradient-corrected
exchange--correlation functionals,''
\emph{Theoretical Chemistry Accounts} \textbf{114}, 145--152 (2005).
\href{https://doi.org/10.1007/s00214-005-0655-y}
{doi:10.1007/s00214-005-0655-y}.

\bibitem{Hamann1979}
D.~R. Hamann, M.~Schl\"uter, and C.~Chiang,
``Norm-conserving pseudopotentials,''
\emph{Physical Review Letters} \textbf{43}, 1494--1497 (1979).
\href{https://doi.org/10.1103/PhysRevLett.43.1494}
{doi:10.1103/PhysRevLett.43.1494}.

\bibitem{Kleinman1982}
L.~Kleinman and D.~M. Bylander,
``Efficacious form for model pseudopotentials,''
\emph{Physical Review Letters} \textbf{48}, 1425--1428 (1982).
\href{https://doi.org/10.1103/PhysRevLett.48.1425}
{doi:10.1103/PhysRevLett.48.1425}.

\bibitem{Payne1992}
M.~C. Payne, M.~P. Teter, D.~C. Allan, T.~A. Arias, and J.~D. Joannopoulos,
``Iterative minimization techniques for ab initio total-energy
calculations: molecular dynamics and conjugate gradients,''
\emph{Reviews of Modern Physics} \textbf{64}, 1045--1097 (1992).
\href{https://doi.org/10.1103/RevModPhys.64.1045}
{doi:10.1103/RevModPhys.64.1045}.

\bibitem{CP2Kmanual}
CP2K developers, ``Pseudopotentials,'' CP2K documentation.
\href{https://manual.cp2k.org/trunk/methods/dft/pseudopotentials.html}
{\nolinkurl{manual.cp2k.org/trunk/methods/dft/pseudopotentials.html}}
(accessed July 2026).

\end{thebibliography}

\end{document}
